\section{Planteamiento y justificación del problema}

\subsection{Planteamiento del problema}
En el ecosistema actual de desarrollo de software, la arquitectura ha dejado de ser un plano estático para convertirse en una entidad dinámica que evoluciona en cada iteración. Ese dinamismo da origen a la deriva arquitectónica, entendida como la divergencia progresiva entre la arquitectura prescrita y la arquitectura realmente implementada en el código fuente \parencite{SanMartin}. En las PyMES de software este fenómeno es especialmente crítico porque el conocimiento arquitectónico suele permanecer distribuido entre código, configuraciones y decisiones tácitas del equipo, en lugar de estar consolidado en documentación viva y verificable \parencite{Hamed2023, Anthony2024, Manglaras2026}.

La literatura reporta que, para 2021, cerca del 65\% del presupuesto de TI de una organización se destinaba al mantenimiento de software \parencite{Pigazzini2021}. Llevado a una referencia ilustrativa de COP 100 millones anuales, esto equivale a COP 65 millones usados solo para sostener el sistema existente. Además, la erosión arquitectónica puede incrementar el esfuerzo y el costo de desarrollo entre un 20\% y un 40\%; en esa misma referencia, el sobrecosto potencial sería de COP 20 a COP 40 millones adicionales \parencite{Pigazzini2021, Jolak2025}. En otras palabras, el problema no es abstracto: la pérdida de integridad arquitectónica se traduce en dinero, tiempo de mantenimiento y menor capacidad de cambio.

Cuando se alude a conceptos técnicos como dependencias cíclicas, componentes Dios o violaciones de capas, en realidad se está nombrando problemas muy concretos. Una dependencia cíclica aparece cuando dos o más módulos terminan dependiendo entre sí y ya no pueden cambiarse o probarse de forma aislada. Un componente Dios es un módulo que concentra demasiadas responsabilidades y demasiadas dependencias salientes, por lo que cualquier cambio en él afecta muchas partes del sistema. Una violación de capas ocurre cuando un módulo rompe la separación esperada entre presentación, negocio y persistencia, debilitando la modularidad y el control del diseño \parencite{SasAvgeriou, Klein2022}. Estos hallazgos son problemáticos porque dificultan las pruebas, aumentan el tiempo de comprensión del sistema, multiplican los efectos colaterales y vuelven más costosa la refactorización.

Sin mecanismos automatizados de vigilancia, estas desviaciones permanecen ocultas hasta que el sistema ya exhibe lentitud para cambiar, fragilidad en el mantenimiento o dependencia excesiva del conocimiento informal del equipo. Por tanto, el problema de investigación no consiste solo en detectar \textit{smells} arquitectónicos aislados, sino en determinar si es posible establecer una gobernanza continua basada exclusivamente en el código fuente, con reglas explícitas y umbrales operativos que permitan distinguir cuándo una arquitectura se mantiene en un estado aceptable y cuándo requiere intervención \parencite{OkooneAgileArchitecture, Klein2022}.

En ese contexto, la pregunta de investigación que orienta este anteproyecto es la siguiente: \textit{\textquestiondown es posible diseñar y evaluar un mecanismo automatizado, basado en análisis estático, grafos de dependencia y \textit{Fitness Functions}, que permita detectar tempranamente derivas arquitectónicas y clasificar el estado de conformidad estructural de un sistema de software para PyMES sin depender de documentación arquitectónica externa?}

\subsection{Justificación del proyecto}
En ausencia de un mecanismo automatizado de detección de deriva arquitectónica, las PyMES de software continúan operando bajo un modelo reactivo de gobernanza técnica, donde la arquitectura solo se revisa cuando ya existen fallos visibles en el sistema. Esta falta de visibilidad estructural genera una acumulación progresiva de desviaciones entre la arquitectura diseñada y la implementada, fenómeno que no siempre se manifiesta primero como error funcional, sino como deterioro silencioso de la mantenibilidad, de la capacidad de evolución y del costo de introducir cambios \parencite{Jolak2025}.

La importancia del proyecto radica en que propone intervenir antes de que ese deterioro se convierta en una trampa de legado. Su aporte no es solamente listar hallazgos técnicos, sino convertir decisiones arquitectónicas que hoy viven de forma tácita en verificaciones automáticas, repetibles y comprensibles para el equipo. Esto es valioso para las PyMES porque reduce la dependencia de documentación externa, disminuye la revisión manual y ofrece una base objetiva para priorizar mantenimiento y refactorización con recursos limitados \parencite{Hamed2023, Klein2022}.

Lo novedoso del trabajo está en la combinación de cuatro elementos dentro de un mismo mecanismo: extracción automática de dependencias desde el código fuente, representación estructural mediante grafos, traducción de decisiones de diseño a \textit{Fitness Functions} y evaluación del estado arquitectónico con un umbral explícito de conformidad. El núcleo del proyecto, por tanto, no es solo detectar desviaciones, sino demostrar que una PyME puede contar con una señal temprana, objetiva y accionable sobre la salud estructural de su software sin depender de reconstrucciones manuales extensas.

Los beneficiarios directos del proyecto serán los equipos de desarrollo, arquitectos de software y líderes técnicos que necesiten evidencia objetiva para identificar desviaciones estructurales y priorizar refactorización. Como beneficiarios indirectos se encuentran las PyMES de software que requieren mecanismos de gobernanza arquitectónica de bajo costo de adopción, con menor dependencia de documentación externa y mayor capacidad para sostener la evolución de sus sistemas en el tiempo. En consecuencia, la utilidad del proyecto no se limita a un artefacto técnico, sino que se proyecta como apoyo a decisiones de mantenimiento, reducción de riesgo estructural y mejora de la sostenibilidad del producto software.

\subsubsection{Hipótesis de trabajo}
Si se implementa un mecanismo automatizado que extraiga dependencias desde el código fuente, ejecute \textit{Fitness Functions} sobre un grafo estructural y clasifique los resultados mediante umbrales explícitos de conformidad, entonces será posible detectar tempranamente derivas arquitectónicas relevantes en un proyecto de software de PyME, diferenciar de manera objetiva un estado aceptable de uno problemático y producir evidencia suficiente para apoyar decisiones de mantenimiento y refactorización sin depender de documentación arquitectónica externa.

